\section{Mitigation Findings}
\label{sec:mitigation}

The detection results raise a natural question: if attention scores are unreliable proxies for grounding, do mitigation interventions reliably reduce hallucination? We evaluate this directly by comparing vanilla generation with three training-free attention interventions and a controlled VCD-greedy decoding baseline. Each method is evaluated on identical POPE samples, with CHAIR captioning additionally reported for the attention-intervention ports, and all deltas are computed relative to the matched vanilla run.
The reproduction scope is deliberately component-level: PAI denotes only its image-attention logit boost, excluding its CFG/contrastive branch; ClearSight denotes VAF mid-layer scaling; VisAttnSink denotes sink redistribution; and VCD-greedy denotes the original/noisy-image contrastive logit rule with greedy decoding rather than stochastic sampling. This scope keeps the test aligned with our question: whether changing visual routing or contrastive decoding verifies object presence.

\begin{table*}[t]
\centering
\small
\setlength{\tabcolsep}{4.2pt}
\begin{tabular}{lrrrrrrrr}
\toprule
Method & $\Delta$Acc & $\Delta$MCC & $\Delta$Yes & $\Delta$TPR & $\Delta$FPR & $\Delta C_i$ & $\Delta$Obj & $\Delta$Hall \\
\midrule
PAI-attn-only & -0.001 & -0.002 & -0.004 & -0.005 & -0.003 & +0.000 & +0.14 & +0.02 \\
ClearSight    & -0.001 & -0.006 & +0.036 & +0.035 & +0.038 & +0.001 & -0.11 & -0.01 \\
VisAttnSink   & -0.004 & -0.008 & +0.006 & +0.002 & +0.009 & +0.009 & +0.24 & +0.10 \\
VCD-greedy    & -0.005 & -0.012 & +0.008 & +0.003 & +0.013 & -- & -- & -- \\
NoLan-compatible & -0.003 & +0.001 & -0.032 & -0.034 & -0.029 & -- & -- & -- \\
\bottomrule
\end{tabular}
\caption{Mitigation behavior relative to vanilla. POPE deltas are macro-averaged over random/popular/adversarial splits; CHAIR deltas are over $5{,}000$ captions where available. ClearSight and VCD-greedy raise recall only by increasing false positives more, while VisAttnSink increases hallucinated object mentions. NoLan-compatible lowers false positives but also lowers recall and the yes rate.}
\label{tab:pope-mitigation}
\end{table*}



\begin{table}[t]
\centering
\scriptsize
\setlength{\tabcolsep}{3.1pt}
\begin{tabular}{lrrrr}
\toprule
Method & $\Delta V$ & $\Delta$TPR & $\Delta$FPR & TP/FP $V$ \\
\midrule
PAI & +0.038 & +0.000 & +0.000 & .103/.105 \\
ClearSight & -0.017 & +0.017 & +0.033 & .063/.063 \\
VisSink & +0.002 & +0.017 & +0.017 & .067/.071 \\
\bottomrule
\end{tabular}
\caption{Adversarial POPE attention-shift audit ($120$ samples). $\Delta V$ is matched active-layer visual-mass change vs.\ vanilla; TP/FP $V$ reports post-intervention visual mass for true and false positives. Visual mass is not selective for correct object presence.}
\label{tab:attention-audit}
\end{table}



\paragraph{POPE gains can be answer-prior shifts.} Table~\ref{tab:pope-mitigation} shows that ClearSight is the only method with a positive macro F1 delta, but its gain follows the pattern of a yes-prior shift. It raises TPR by $3.5$ points and FPR by $3.8$ points, while MCC decreases. On adversarial POPE specifically, ClearSight increases TPR by $3.5$ points but FPR by $5.5$ points, causing accuracy to drop by one point. Thus the method makes the model more likely to answer ``yes,'' but does not improve object-presence discrimination. PAI-attention-only is nearly neutral but slightly worse, suggesting that PAI's full reported benefit may depend on its contrastive/CFG component rather than attention amplification alone. VisAttnSink mildly increases yes rate and FPR while reducing accuracy and MCC.

\paragraph{Captioning does not show hidden mitigation.} The CHAIR columns of Table~\ref{tab:pope-mitigation} rule out the possibility that POPE misses a captioning improvement. None of the interventions lowers CHAIR$_i$. PAI-attention-only and VisAttnSink make captions more object-rich; the latter increases hallucinated object mentions by $0.10$ per caption on average. ClearSight slightly shortens captions and mentions fewer objects, but still does not reduce CHAIR. These patterns support the same interpretation as the detection study: visual attention manipulation changes output behavior, but does not provide calibrated evidence verification for object claims.

\paragraph{Attention changes are real but not selective.} Table~\ref{tab:attention-audit} shows why the behavioral gains are misleading. PAI increases matched active-layer visual mass by $+0.038$ and reduces prefix/sink mass, yet every adversarial POPE decision is unchanged. Moreover, PAI false positives receive slightly \emph{higher} visual mass than true positives ($0.105$ vs.\ $0.103$). ClearSight raises TPR but raises FPR twice as much on the same subset, reducing MCC; its TP and FP visual masses are nearly identical. Visual Attention Sink moves attention in the intended direction only weakly ($+0.002$ visual mass, $-0.005$ sink mass) and increases TPR and FPR by the same amount. The interventions therefore change routing, but the routing change is not selective for object presence.

\paragraph{Why the interventions fail.} The methods fail for related but distinct reasons. PAI-attention-only adds a positive offset proportional to the magnitude of image-attention logits, so it strengthens whatever image-attention pattern already exists; it does not test whether the attended patch supports the queried object. ClearSight applies a broader mid-layer amplification of text-to-image attention while suppressing prefix attention, which can make image evidence more influential but can also shift the model toward affirmative answers whenever the question mentions a plausible object. Visual Attention Sink redistributes mass from high-activation sink tokens to visual tokens, increasing visual-token participation during decoding, but this can make generation more object-rich rather than more faithful. VCD-greedy contrasts original-image logits with noised-image logits, reducing one form of language-prior reliance, but it still does not ask whether the visual evidence distinguishes the target from associated objects. In all cases, the operation changes routing or decoding pressure without introducing an explicit object-presence verifier. This explains why POPE recall and false-positive rate move together, and why richer captions can contain more hallucinated objects.
